\chapter{Funções de referência}\label{ch:g10:reffunc}

Um punhado de funções simples --- afim, quadrado, inverso, raiz quadrada,
cubo --- aparece em toda parte, sozinho ou combinado. Saber de cor seus
\omterm{def:g10:functions:graph}{gráficos} e sua variação transforma muitos problemas em um esboço rápido.
Este capítulo estuda cada uma delas por sua vez e demonstra sua variação
com a técnica de comparação do \cref{ch:g10:functions}.

\section{Funções afins}

\begin{definition}[Função afim]\label{def:g10:reffunc:affine}
Uma \emph{função afim}\index{função afim} é uma \omterm{def:g10:functions:function}{função} da forma
\[
f(x) = mx + p,
\]
onde $m$ e $p$ são \omterm{def:g10:numbers:sets}{números reais} fixos. Seu \omterm{def:g10:functions:graph}{gráfico} é uma reta: $m$ é o
\emph{coeficiente angular}\index{coeficiente angular} e $p$ o
\emph{coeficiente linear}\index{coeficiente linear} (a reta corta o eixo
vertical em $(0, p)$). Quando $p = 0$, a \omterm{def:g10:functions:function}{função} é \emph{linear}:
$f(x) = mx$ exprime uma proporcionalidade.
\end{definition}

\begin{proposition}[Coeficiente angular e variação]\label{prop:g10:reffunc:affine}
Seja $f(x) = mx + p$.
\begin{enumerate}
  \item Para quaisquer duas entradas distintas $u \neq v$:
        $m = \dfrac{f(v) - f(u)}{v - u}$ (o \omterm{def:g10:reffunc:affine}{coeficiente angular} é a
        variação da saída por unidade de variação da entrada).
  \item Se $m > 0$, $f$ é estritamente \omterm{def:g10:functions:variations}{crescente} em $\R$; se $m < 0$,
        estritamente \omterm{def:g10:functions:variations}{decrescente}; se $m = 0$, constante.
\end{enumerate}
\end{proposition}

\begin{proof}
1. Calcule $f(v) - f(u) = (mv + p) - (mu + p) = m(v - u)$ e divida em
seguida por $v - u \neq 0$.

2. Se $u < v$, então $f(v) - f(u) = m(v-u)$ tem o sinal de $m$, pois
$v - u > 0$: com $m$ positivo, $f(u) < f(v)$ (\omterm{def:g10:functions:variations}{crescente}); com $m$
negativo, $f(u) > f(v)$ (\omterm{def:g10:functions:variations}{decrescente}).
\end{proof}

\begin{omfigure}
\begin{tikzpicture}
\begin{axis}[omaxis,
  width=7.8cm, height=5.6cm,
  xmin=-3.4, xmax=3.4, ymin=-2.8, ymax=4.4,
  xtick={-2,-1,1,2}, ytick={-1,1,2,3}]
  \addplot[omDef, thick, domain=-3:2.6] {0.5*x + 1};
  \addplot[omThm, thick, domain=-2.2:3] {-x + 1};
  \node[omDef, font=\small] at (axis cs:2.5,2.7) {$y = \tfrac12 x + 1$};
  \node[omThm, font=\small] at (axis cs:-2.0,2.6) {$y = -x + 1$};
  \draw[dashed] (axis cs:1,1.5) -- (axis cs:2,1.5) -- (axis cs:2,2);
  \node[font=\footnotesize, below] at (axis cs:1.5,1.5) {$1$};
  \node[font=\footnotesize, right] at (axis cs:2,1.75) {$m$};
\end{axis}
\end{tikzpicture}

{\small Duas \omterm{def:g10:reffunc:affine}{funções afins}: \omterm{def:g10:functions:variations}{crescente} para $m = \frac12 > 0$, \omterm{def:g10:functions:variations}{decrescente}
para $m = -1 < 0$. Avançar uma unidade para a direita faz $y$ variar de
$m$.}
\end{omfigure}

\begin{example}\label{ex:g10:reffunc:affine}
Encontre a \omterm{def:g10:reffunc:affine}{função afim} cujo \omterm{def:g10:functions:graph}{gráfico} passa por $A(1, 3)$ e $B(4, 9)$. O
\omterm{def:g10:reffunc:affine}{coeficiente angular} é
\[
m = \frac{9 - 3}{4 - 1} = \frac{6}{3} = 2 .
\]
Então $f(x) = 2x + p$, e $f(1) = 3$ dá $2 + p = 3$, logo $p = 1$:
$f(x) = 2x + 1$. Confira com $B$: $f(4) = 9$.
\end{example}

\section{A função quadrado}

\begin{proposition}[Função quadrado]\label{prop:g10:reffunc:square}
A \omterm{def:g10:functions:function}{função} $f(x) = x^2$, definida em $\R$:
\begin{enumerate}
  \item é estritamente \omterm{def:g10:functions:variations}{decrescente} em $\intoc{-\infty}{0}$ e estritamente
        \omterm{def:g10:functions:variations}{crescente} em $\intco{0}{+\infty}$, com \omterm{def:g10:functions:extrema}{mínimo} $0$ em $x = 0$;
  \item satisfaz $f(-x) = f(x)$ para todo $x$: seu \omterm{def:g10:functions:graph}{gráfico}, uma
        \emph{parábola}\index{parábola}, é simétrico em relação ao eixo
        vertical.
\end{enumerate}
\end{proposition}

\begin{proof}
1. Sejam $0 \leq u < v$. Então
\[
f(v) - f(u) = v^2 - u^2 = (v - u)(v + u) > 0,
\]
pois $v - u > 0$ e $v + u > 0$: $f$ é estritamente \omterm{def:g10:functions:variations}{crescente} em
$\intco{0}{+\infty}$. Se $u < v \leq 0$, então $v - u > 0$, mas
$v + u < 0$, logo $f(v) - f(u) < 0$: estritamente \omterm{def:g10:functions:variations}{decrescente}. Por fim,
$x^2 \geq 0 = f(0)$ para todo $x$.

2. $(-x)^2 = x^2$; assim, os pontos $(x, x^2)$ e $(-x, x^2)$, simétricos
em relação ao eixo vertical, estão ambos no \omterm{def:g10:functions:graph}{gráfico}.
\end{proof}

\begin{remark}[Quadrados e ordem]
Como a \omterm{def:g10:functions:function}{função} quadrado decresce no lado negativo, elevar ao quadrado
\emph{inverte} a ordem dos números negativos: $-3 < -2$, mas $9 > 4$.
Nunca eleve os dois lados de uma desigualdade ao quadrado sem antes
verificar os sinais.
\end{remark}

\section{A função inverso}

\begin{proposition}[Função inverso]\label{prop:g10:reffunc:inverse}
A \omterm{def:g10:functions:function}{função} $f(x) = \dfrac1x$, definida para $x \neq 0$:
\begin{enumerate}
  \item é estritamente \omterm{def:g10:functions:variations}{decrescente} em $\intoo{-\infty}{0}$ e
        estritamente \omterm{def:g10:functions:variations}{decrescente} em $\intoo{0}{+\infty}$;
  \item satisfaz $f(-x) = -f(x)$: seu \omterm{def:g10:functions:graph}{gráfico}, uma
        \emph{hipérbole}\index{hipérbole}, é simétrico em relação à
        origem.
\end{enumerate}
\end{proposition}

\begin{proof}
1. Sejam $0 < u < v$. Então
\[
f(v) - f(u) = \frac1v - \frac1u = \frac{u - v}{uv} < 0,
\]
pois $u - v < 0$ e $uv > 0$: estritamente \omterm{def:g10:functions:variations}{decrescente} em
$\intoo{0}{+\infty}$. Em $\intoo{-\infty}{0}$ o mesmo quociente tem
novamente $u - v < 0$ e $uv > 0$ (produto de dois negativos), de modo que
$f$ também decresce ali.

2. $\frac{1}{-x} = -\frac1x$.
\end{proof}

\begin{remark}
Cuidado: $\frac1x$ \emph{não} é \omterm{def:g10:functions:variations}{decrescente} em todo o seu \omterm{def:g10:functions:function}{domínio}. De
$u = -1$ até $v = 1$ o valor salta de $-1$ para $1$. Os dois ramos devem
ser estudados separadamente.
\end{remark}

\begin{omfigure}
\begin{tikzpicture}
\begin{axis}[omaxis,
  width=5.6cm, height=5.2cm,
  xmin=-3.2, xmax=3.2, ymin=-2.6, ymax=6.4,
  xtick={-2,-1,1,2}, ytick={1,4},
  title={$y = x^2$}]
  \addplot[omDef, thick, domain=-2.45:2.45, samples=90] {x^2};
\end{axis}
\end{tikzpicture}\hfill
\begin{tikzpicture}
\begin{axis}[omaxis,
  width=5.6cm, height=5.2cm,
  xmin=-3.4, xmax=3.4, ymin=-3.4, ymax=3.4,
  xtick={-1,1}, ytick={-1,1},
  title={$y = \frac1x$}]
  \addplot[omDef, thick, domain=0.3:3.1, samples=80] {1/x};
  \addplot[omDef, thick, domain=-3.1:-0.3, samples=80] {1/x};
\end{axis}
\end{tikzpicture}

{\small A parábola $y = x^2$ (simétrica em relação ao eixo vertical) e a
hipérbole $y = \frac1x$ (dois ramos, simétricos em relação à origem).}
\end{omfigure}

\section{Raiz quadrada e cubo}

\begin{proposition}[Função raiz quadrada]\label{prop:g10:reffunc:sqrt}
A \omterm{def:g10:functions:function}{função} $f(x) = \sqrt{x}$, definida em $\intco{0}{+\infty}$, é
estritamente \omterm{def:g10:functions:variations}{crescente}.
\end{proposition}

\begin{proof}
Sejam $0 \leq u < v$. Multiplique e divida pelo \emph{conjugado}:
\[
\sqrt v - \sqrt u
= \frac{(\sqrt v - \sqrt u)(\sqrt v + \sqrt u)}{\sqrt v + \sqrt u}
= \frac{v - u}{\sqrt v + \sqrt u} > 0,
\]
pois $v - u > 0$ e $\sqrt v + \sqrt u > 0$ (note que $v > 0$, logo
$\sqrt v > 0$).
\end{proof}

\begin{proposition}[Função cubo]\label{prop:g10:reffunc:cube}
A \omterm{def:g10:functions:function}{função} $f(x) = x^3$, definida em $\R$, é estritamente \omterm{def:g10:functions:variations}{crescente}, e seu
\omterm{def:g10:functions:graph}{gráfico} é simétrico em relação à origem.
\end{proposition}
\admitted

\begin{omfigure}
\begin{tikzpicture}
\begin{axis}[omaxis,
  width=5.6cm, height=5.0cm,
  xmin=-0.7, xmax=5.4, ymin=-0.7, ymax=2.9,
  xtick={1,4}, ytick={1,2},
  title={$y = \sqrt x$}]
  \addplot[omDef, thick, domain=0:5.1, samples=100] {sqrt(x)};
\end{axis}
\end{tikzpicture}\hfill
\begin{tikzpicture}
\begin{axis}[omaxis,
  width=5.6cm, height=5.0cm,
  xmin=-2.2, xmax=2.2, ymin=-3.4, ymax=3.4,
  xtick={-1,1}, ytick={-1,1},
  title={$y = x^3$}]
  \addplot[omDef, thick, domain=-1.48:1.48, samples=90] {x^3};
\end{axis}
\end{tikzpicture}

{\small A raiz quadrada (definida só para $x \geq 0$) e o cubo, ambas
\omterm{def:g10:functions:variations}{crescentes}.}
\end{omfigure}

\begin{proposition}[Comparação entre $x$, $x^2$ e $\sqrt x$]\label{prop:g10:reffunc:compare}
Para $0 < x < 1$:\ \ $x^2 < x < \sqrt x$.\quad
Para $x > 1$:\ \ $\sqrt x < x < x^2$.\quad
Em $x = 0$ e $x = 1$ os três valores coincidem.
\end{proposition}

\begin{proof}
Suponha $0 < x < 1$. Multiplicar $x < 1$ por $x > 0$ dá $x^2 < x$. Em
seguida, $x < \sqrt x$: como os dois lados são positivos e elevar ao
quadrado é \omterm{def:g10:functions:variations}{crescente} nos positivos, essa desigualdade equivale a
$x^2 < x$, que acabamos de provar. Para $x > 1$, multiplicar $x > 1$ por
$x$ dá $x^2 > x$, e $\sqrt x < x$ segue pelo mesmo argumento com
quadrados.
\end{proof}

\begin{omfigure}
\begin{tikzpicture}
\begin{axis}[omaxis,
  width=7.4cm, height=5.8cm,
  xmin=-0.35, xmax=2.5, ymin=-0.35, ymax=2.9,
  xtick={1,2}, ytick={1,2}]
  \addplot[omThm, thick, domain=0:1.65, samples=80] {x^2};
  \addplot[omDef, thick, domain=0:2.35] {x};
  \addplot[omMeth, thick, domain=0:2.35, samples=80] {sqrt(x)};
  \fill (axis cs:1,1) circle (1.5pt);
  \node[omThm, font=\small, right] at (axis cs:1.45,2.4) {$y = x^2$};
  \node[omDef, font=\small, right] at (axis cs:2.0,1.85) {$y = x$};
  \node[omMeth, font=\small, right] at (axis cs:2.0,1.28) {$y = \sqrt x$};
\end{axis}
\end{tikzpicture}

{\small As três curvas se cruzam em $(0,0)$ e $(1,1)$ e trocam de ordem
nesses pontos: entre $0$ e $1$, $x^2$ é a menor e $\sqrt x$ a maior;
depois de $1$ a ordem se inverte.}
\end{omfigure}

\begin{method}[Comparar imagens]\label{met:g10:reffunc:compare}
Para comparar $f(a)$ e $f(b)$ para uma \omterm{def:g10:functions:function}{função} de referência $f$:
\begin{enumerate}
  \item coloque $a$ e $b$ em um \omterm{def:g10:numbers:interval}{intervalo} onde a variação de $f$ seja
        conhecida;
  \item se $f$ cresce ali, as \omterm{def:g10:functions:function}{imagens} ficam na mesma ordem que as
        entradas; se $f$ decresce, a ordem se inverte;
  \item se $a$ e $b$ não estiverem no mesmo \omterm{def:g10:numbers:interval}{intervalo} de monotonicidade,
        trate primeiro os sinais (por exemplo, no quadrado: uma entrada
        negativa e uma positiva).
\end{enumerate}
\end{method}

\begin{example}
Compare $\dfrac{1}{2.3}$ e $\dfrac{1}{2.4}$: as duas entradas estão em
$\intoo{0}{+\infty}$, onde a \omterm{def:g10:functions:function}{função} inverso decresce, e $2.3 < 2.4$, logo
$\dfrac{1}{2.3} > \dfrac{1}{2.4}$. Compare $(-1.2)^2$ e $(-1.5)^2$: em
$\intoc{-\infty}{0}$ o quadrado decresce, e $-1.5 < -1.2$, logo
$(-1.5)^2 > (-1.2)^2$ --- de fato, $2.25 > 1.44$.
\end{example}

\section{Exercícios}

\begin{exercise}[$\star$]\label{exo:g10:reffunc:1}
Para cada \omterm{def:g10:reffunc:affine}{função afim}, dê o \omterm{def:g10:reffunc:affine}{coeficiente angular}, o \omterm{def:g10:reffunc:affine}{coeficiente linear} e
diga se ela é \omterm{def:g10:functions:variations}{crescente} ou \omterm{def:g10:functions:variations}{decrescente}:
\[
f(x) = 3x - 2, \qquad
g(x) = -\tfrac12 x + 4, \qquad
h(x) = 7, \qquad
k(x) = -x .
\]
\end{exercise}

\begin{exercise}[$\star$]\label{exo:g10:reffunc:2}
Encontre a \omterm{def:g10:reffunc:affine}{função afim} cujo \omterm{def:g10:functions:graph}{gráfico} passa por $(2, 1)$ e $(5, 10)$; e
depois a que passa por $(-1, 4)$ e $(3, -4)$.
\end{exercise}

\begin{exercise}[$\star$]\label{exo:g10:reffunc:3}
Sem calculadora, compare:
\[
(3.1)^2 \text{ e } (3.2)^2; \qquad
(-2.7)^2 \text{ e } (-2.8)^2; \qquad
\frac{1}{5.1} \text{ e } \frac{1}{5.2}; \qquad
\sqrt{17} \text{ e } \sqrt{15}.
\]
\end{exercise}

\begin{exercise}[$\star$]\label{exo:g10:reffunc:4}
Usando o \omterm{def:g10:functions:graph}{gráfico} da \omterm{def:g10:functions:function}{função} quadrado, resolva $x^2 = 16$, depois
$x^2 \leq 16$ e depois $x^2 > 9$.
\end{exercise}

\begin{exercise}[$\star$]\label{exo:g10:reffunc:5}
Seja $x \in \intoo{0}{1}$. Ordene os números $x$, $x^2$, $x^3$ e
$\sqrt x$ do menor para o maior e confira sua resposta com $x = 0.25$.
\end{exercise}

\begin{exercise}[$\star\star$]\label{exo:g10:reffunc:6}
Resolva graficamente e depois algebricamente: $\dfrac1x = 2$ e
$\dfrac1x < 2$ para $x > 0$. O que muda se também permitirmos $x < 0$?
\end{exercise}

\begin{exercise}[$\star\star$]\label{exo:g10:reffunc:7}
Sabendo que a \omterm{def:g10:functions:function}{função} quadrado decresce em $\intoc{-\infty}{0}$ e cresce
em $\intco{0}{+\infty}$, enquadre $x^2$ quando:
\[
\text{(a) } x \in \intcc{2}{5};
\qquad
\text{(b) } x \in \intcc{-3}{-1};
\qquad
\text{(c) } x \in \intcc{-2}{3}.
\]
(No caso (c), atenção: o \omterm{def:g10:functions:extrema}{mínimo} de $x^2$ não está em um extremo.)
\end{exercise}

\begin{exercise}[$\star\star$]\label{exo:g10:reffunc:8}
Uma empresa de táxi cobra uma taxa fixa de $4$ mais $1.5$ por
quilômetro; outra não cobra taxa e cobra $2$ por quilômetro. Modele cada
preço por uma \omterm{def:g10:reffunc:affine}{função afim} da distância $x$, trace as duas retas e
descubra a partir de que distância a primeira empresa sai mais barata.
\end{exercise}

\begin{exercise}[$\star\star$]\label{exo:g10:reffunc:9}
Mostre que, para todos $a, b \geq 0$: $\sqrt{a + b} \leq \sqrt a +
\sqrt b$. (Sugestão: os dois lados são não negativos, portanto compare
seus quadrados.) Quando há igualdade?
\end{exercise}

\begin{exercise}[$\star\star\star$]\label{exo:g10:reffunc:10}
Seja $f(x) = \dfrac{2x + 1}{x - 1}$, definida para $x \neq 1$.
\begin{enumerate}
  \item Mostre que $f(x) = 2 + \dfrac{3}{x - 1}$ para todo $x \neq 1$.
  \item Deduza a variação de $f$ em $\intoo{1}{+\infty}$ a partir da
        variação da \omterm{def:g10:functions:function}{função} inverso.
\end{enumerate}
\end{exercise}

\section{Problema: As leis da natureza falam em funções de referência}

\begin{problem}[{Problema de fim de semana --- distâncias de frenagem
crescem como $v^2$, alavancas obedecem a $1/d$, planetas seguem
$a^{3/2}$: ler a física com as funções de referência}]
\label{pb:g10:reffunc:1}
Abra um livro de física e os mesmos poucos \omterm{def:g10:functions:graph}{gráficos} reaparecem em
cada página: a parábola, a hipérbole, a curva da raiz. A natureza
escreve suas leis com as funções de referência deste capítulo ---
e conhecer suas formas (\cref{prop:g10:reffunc:square},
\cref{prop:g10:reffunc:inverse}, \cref{prop:g10:reffunc:sqrt},
\cref{prop:g10:reffunc:cube}) basta para frear um carro,
equilibrar uma alavanca e cronometrar os planetas. O grande final
é a terceira lei de Kepler, lida diretamente nos dados do sistema
solar.

\medskip
\noindent\textbf{Parte I --- A lei quadrática da frenagem.}
Uma regra prática para a distância de frenagem de um carro em
pista seca: $d = \left(\frac{v}{10}\right)^{2}$ metros, à
velocidade $v$ km/h.
\begin{enumerate}
  \item Calcule as distâncias de frenagem a $30$, $50$, $90$ e
        $130$ km/h.
  \item Dobrar a velocidade multiplica a distância de frenagem
        por quanto? Explique a partir da escala da \omterm{def:g10:functions:function}{função}
        quadrado e confira em $50 \to 100$ km/h.
  \item Perícia: as marcas de derrapagem medem $50$ m. De que
        velocidade a fórmula acusa o motorista (até o km/h)? Qual
        \omterm{def:g10:functions:function}{função} de referência respondeu --- e em que \omterm{def:g10:functions:function}{domínio} a
        leitura é inequívoca
        (\cref{prop:g10:reffunc:sqrt})?
  \item Compare a distância extra causada por \emph{um} km/h a
        mais, em baixa e em alta velocidade: calcule
        $d(31) - d(30)$ e $d(91) - d(90)$. Que propriedade da
        parábola (\cref{prop:g10:reffunc:square}) as duas
        respostas ilustram?
  \item A parada real acrescenta o tempo de reação --- cerca de
        um segundo, durante o qual o carro percorre $0.28v$
        metros: $D(v) = 0.28v + \left(\frac{v}{10}\right)^2$.
        Calcule $D(50)$ e $D(130)$. Qual dos dois termos --- o
        afim ou o quadrático --- domina na velocidade urbana, e
        qual domina na rodovia?
\end{enumerate}

\medskip
\noindent\textbf{Parte II --- As hipérboles do dia a dia.}
\begin{enumerate}[resume]
  \item Uma alavanca se equilibra quando força $\times$ distância
        é a mesma dos dois lados. Uma criança de $60$ kg senta a
        $2$ m do apoio; a força equilibrante à distância $d$ do
        outro lado é $F(d) = \frac{120}{d}$ (em equivalentes de
        quilograma). Calcule $F(0.5)$, $F(1)$, $F(3)$.
  \item O que a variação da \omterm{def:g10:functions:function}{função} inverso
        (\cref{prop:g10:reffunc:inverse}) diz sobre as
        alavancas --- e por que a bravata de Arquimedes
        (``dai-me um ponto de apoio e moverei a Terra'')
        esconde-se na cauda da hipérbole? O que proíbe $d = 0$?
  \item O som e a luz enfraquecem com o \emph{quadrado} da
        distância: a $1$ m de uma lâmpada a intensidade é $100$
        unidades; dê-a a $2$, $3$ e $10$ m. Explique o expoente
        $2$ com uma esfera: sobre que área a luz se espalhou à
        distância $d$ (o problema sobre a lápide de Arquimedes,
        no volume do ensino fundamental)?
  \item Um ônibus para a excursão escolar custa $600$ euros,
        divididos igualmente entre $n$ participantes. Faça uma
        tabela do custo por pessoa para $n = 10, 20, 30$,
        descubra quantos participantes o reduzem a $25$ euros ou
        menos e diga o que a forma da hipérbole promete --- e
        recusa --- à medida que $n$ cresce.
  \item Produzir $x$ cartazes custa $2x + 3$ euros ($3$ euros de
        preparação). Mostre que o custo \emph{médio} por cartaz é
        $a(x) = 2 + \frac3x$, descreva sua variação e interprete
        economicamente sua assíntota horizontal. (Compare com a
        álgebra do \cref{exo:g10:reffunc:10}.)
\end{enumerate}

\medskip
\noindent\textbf{Parte III --- A harmonia de Kepler.}
Para cada planeta, seja $a$ sua distância média ao Sol (em
unidades astronômicas, Terra $= 1$) e $T$ seu período (em anos).
Dados: Marte $a = 1.52$, $T = 1.88$; Júpiter $a = 5.20$,
$T = 11.86$; Saturno $a = 9.54$, $T = 29.4$.
\begin{enumerate}[resume]
  \item Calcule $a^3$ e $T^2$ para a Terra, Marte e Júpiter. O
        que Kepler notou em 1618?
  \item Enuncie a lei como \omterm{def:g10:functions:function}{função}: $T = \sqrt{a^3} = a\sqrt a$.
        Teste-a em Saturno.
  \item Duas previsões: um asteroide orbita a $a = 4$ UA ---
        encontre seu período (os números saem inteiros); um
        cometa retorna a cada $27$ anos --- encontre sua
        distância média (procure um cubo perfeito).
  \item Que três funções de referência a lei entrelaça? E sua
        regra de escala: quando $a$ é multiplicado por $4$, o que
        acontece com $T$? (Confira com o asteroide, comparado à
        Terra.)
  \item Kepler leu essa lei nos dados de Tycho Brahe setenta anos
        antes de a gravitação de Newton explicá-la. Em uma frase:
        o que esse episódio diz sobre o poder de reconhecer uma
        \omterm{def:g10:functions:function}{função} de referência em uma tabela de números?
\end{enumerate}

\medskip
\noindent\textbf{Parte IV --- A tartaruga e a lebre.}
\begin{enumerate}[resume]
  \item Ordene $\sqrt x$, $x$ e $x^2$ em $\intoo{0}{1}$ e em
        $\intoo{1}{+\infty}$
        (\cref{prop:g10:reffunc:compare}) e verifique as duas
        ordens em $x = 0.25$ e $x = 4$.
  \item Resolva $\sqrt x > x$ por completo (para $x \geq 0$,
        eleve ao quadrado com cuidado e conclua com um \omterm{def:g10:algebra:signtable}{quadro de
        sinais}).
  \item Traga o cubo: ordene $x$, $x^2$, $x^3$ em $x = 0.5$ e em
        $x = 2$, e encontre todos os pontos onde as funções
        quadrado e cubo se cruzam.
  \item Dois planos de poupança pagam, após $t$ anos,
        $100\sqrt t$ euros (plano A) ou $10 t^2$ euros (plano B).
        Mostre que B ultrapassa A quando $t^3 = 100$ e dê o
        instante da ultrapassagem com precisão de meses. Moral
        sobre raízes contra potências no longo prazo?
  \item Final: para cada lei deste problema --- frenagem
        ($v^2$), alavanca ($1/d$), lâmpada ($1/d^2$), Kepler
        ($a^{3/2}$), plano A ($\sqrt t$) --- diga em uma oração
        qual traço do seu \omterm{def:g10:functions:graph}{gráfico} de referência carrega o
        significado físico (subida cada vez mais íngreme,
        assíntota vertical, esfera que se espalha, expoente de
        escala, crescimento que se achata). Conclusão: as funções
        de referência são o alfabeto; a natureza escreve com
        elas.
\end{enumerate}
\end{problem}
